
\documentclass{article}
\usepackage{colortbl}
\usepackage{makecell}
\usepackage{multirow}
\usepackage{supertabular}

\begin{document}

\newcounter{utterance}

\centering \large Interaction Transcript for game `cladder', experiment `full\_v1.5\_default', episode 7019 with selfrefine{-}Qwen3.5{-}27B.
\vspace{24pt}

{ \footnotesize  \setcounter{utterance}{1}
\setlength{\tabcolsep}{0pt}
\begin{supertabular}{c@{$\;$}|p{.15\linewidth}@{}p{.15\linewidth}p{.15\linewidth}p{.15\linewidth}p{.15\linewidth}p{.15\linewidth}}
    \# & \multicolumn{2}{c}{Player} && \multicolumn{2}{c}{Game Master} \\
    \hline

    \theutterance \stepcounter{utterance}  
    & & & \multicolumn{4}{p{0.6\linewidth}}{
        \cellcolor[rgb]{0.9,0.9,0.9}{
            \makecell[{{p{\linewidth}}}]{
                \texttt{\tiny{[P1$\langle$GM]}}
                \texttt{You are an expert in causal inference. The following question is not a typical commonsense query, but rather a meticulously designed question created by a professor specializing in causal inference, intended to assess the students' mastery of the course content.} \\
\texttt{Imagine a self{-}contained, hypothetical world with only the following conditions, and without any unmentioned factors or causal relationships: Hwax has a direct effect on gyzp and jyka. Jyka has a direct effect on lirg. Gyzp has a direct effect on lirg. Method 1: We look at how jyka correlates with lirg case by case according to gyzp. Method 2: We look directly at how jyka correlates with lirg in general. To understand how jyka affects lirg, is it more correct to use the Method 1 than Method 2?} \\
\texttt{Start your answer with 'Yes' or 'No', followed by additional reasoning or evidence to support your explanation.} \\
            }
        }
    }
     \\ \\

    \theutterance \stepcounter{utterance}  
    & \multicolumn{4}{p{0.6\linewidth}}{
        \cellcolor[rgb]{0.561,0.737,0.561}{
            \makecell[{{p{\linewidth}}}]{
                \texttt{\tiny{[P1$\rangle$GM]}}
                \texttt{No} \\
\texttt{In the specified causal structure (\$Hwax \textbackslash{}to Gyzp\$, \$Hwax \textbackslash{}to Jyka\$, \$Jyka \textbackslash{}to Lirg\$, \$Gyzp \textbackslash{}to Lirg\$), \$Gyzp\$ is a mediator on the path from \$Hwax\$ to \$Lirg\$, but more critically for the relationship between \$Jyka\$ and \$Lirg\$, \$Gyzp\$ is a **collider** on the path \$Jyka \textbackslash{}leftarrow Hwax \textbackslash{}to Gyzp \textbackslash{}leftarrow (\textbackslash{}text\{no other parent\}) \textbackslash{}to Lirg\$? No, let's trace the paths strictly.} \\
\texttt{The path of interest is \$Jyka \textbackslash{}to Lirg\$.} \\
\texttt{The backdoor path is \$Jyka \textbackslash{}leftarrow Hwax \textbackslash{}to Gyzp \textbackslash{}to Lirg\$.} \\
\texttt{In this backdoor path, \$Gyzp\$ is not a collider; it is an intermediate node on a chain (\$Hwax \textbackslash{}to Gyzp \textbackslash{}to Lirg\$). However, the critical error in Method 1 is that \$Gyzp\$ is a **mediator** of the effect of \$Hwax\$ on \$Lirg\$, but wait, does \$Gyzp\$ mediate the effect of \$Jyka\$? No.} \\
\texttt{Let's re{-}evaluate the role of \$Gyzp\$ relative to the \$Jyka \textbackslash{}to Lirg\$ relationship.} \\
\texttt{The backdoor path is \$Jyka \textbackslash{}leftarrow Hwax \textbackslash{}to Gyzp \textbackslash{}to Lirg\$.} \\
\texttt{To block this backdoor path, we must condition on a variable that blocks the flow. Conditioning on \$Hwax\$ would block it. Conditioning on \$Gyzp\$ also blocks this specific backdoor path because \$Gyzp\$ is a non{-}collider on the path \$Hwax \textbackslash{}to Gyzp \textbackslash{}to Lirg\$.} \\
\texttt{However, is there another path?} \\
\texttt{The graph is: \$Hwax \textbackslash{}to Jyka \textbackslash{}to Lirg\$ and \$Hwax \textbackslash{}to Gyzp \textbackslash{}to Lirg\$.} \\
\texttt{Wait, the prompt says: "Hwax has a direct effect on gyzp and jyka. Jyka has a direct effect on lirg. Gyzp has a direct effect on lirg."} \\
\texttt{So the paths from \$Jyka\$ to \$Lirg\$ are:} \\
\texttt{1. Direct: \$Jyka \textbackslash{}to Lirg\$.} \\
\texttt{2. Backdoor: \$Jyka \textbackslash{}leftarrow Hwax \textbackslash{}to Gyzp \textbackslash{}to Lirg\$.} \\
\texttt{Method 2 (unadjusted) estimates the total association, which includes the direct effect plus the bias from the backdoor path.} \\
\texttt{Method 1 conditions on \$Gyzp\$.} \\
\texttt{Does conditioning on \$Gyzp\$ block the} \\
            }
        }
    }
    & & \\ \\

    \theutterance \stepcounter{utterance}  
    & & & \multicolumn{2}{p{0.3\linewidth}}{
        \cellcolor[rgb]{0.9,0.9,0.9}{
            \makecell[{{p{\linewidth}}}]{
                \texttt{\tiny{[GM$|$GM]}}
                \texttt{No} \\
            }
        }
    }
    & & \\ \\

    \theutterance \stepcounter{utterance}  
    & & & \multicolumn{2}{p{0.3\linewidth}}{
        \cellcolor[rgb]{0.9,0.9,0.9}{
            \makecell[{{p{\linewidth}}}]{
                \texttt{\tiny{[GM$|$GM]}}
                \texttt{yes} \\
            }
        }
    }
    & & \\ \\

    \theutterance \stepcounter{utterance}  
    & & & \multicolumn{2}{p{0.3\linewidth}}{
        \cellcolor[rgb]{0.9,0.9,0.9}{
            \makecell[{{p{\linewidth}}}]{
                \texttt{\tiny{[GM$|$GM]}}
                \texttt{game\_result = LOSE} \\
            }
        }
    }
    & & \\ \\

\end{supertabular}
}

\end{document}
